\section*{Métrica inducida por el espacio ambiente ($\mathbb{R}^n$) en una superficie $\Sigma$ ($\mathbb{R}^3$)}
\underline{Cálculo detallado}:
\begin{align*}
    &1.\;  \text{Parametriza superficie mediante inmersión } X:U\subset \mathbb{R}^2 \longrightarrow \mathbb{R}^n, \; (u,v)\mapsto X(u,v)\\
    &2. \;  \text{Calcula los vectores tangentes: }
    X_u = \frac{\partial X}{\partial u}, \; X_v= \frac{\partial X}{\partial 
    v}\\
    &3.\; \text{Restringe la métrica: } \!E\!=\! \langle X_u,X_u\rangle\;\; F\!=\! \langle X_u,X_v\rangle\; \; G\!=\! \langle X_v,X_v\rangle \;\; \langle\cdot ,\cdot\rangle\!\equiv\!\text{prod.esc.}\\
    &4.\; \text{Tensor métrico inducido: } g=
    \begin{pmatrix}
        E&F\\
        F&G
    \end{pmatrix}\; \; \; \; ds^2 = E\,du^2+ 2F\,du\,dv + G\, dv^2\\
    &5.\;  \text{Elemento de área:} dA= \sqrt{\text{det}\,g} \: du\,dv
\end{align*}
\underline{Ejemplo (toroide)} \text{ radio circunferencia: $a$, distancia origen-centro circunferencia: $b$}
\begin{align*}
    &\text{Parametrización:}\hspace{-5pt}
    &\begin{cases} 
x= (b+a\,\cos\alpha)\cos\theta
\\
y= (b+a\,\cos\alpha)\sin\theta
\\
z= a\sin\alpha
\end{cases}
\left|
\hspace{-12pt}
\begin{array}{ll}
     & \alpha\equiv \text{rotación en circunferencia}\\
     &\, \theta\equiv \text{rotación eje z} 
\end{array}
\right.
\end{align*}
\begin{align*}
X_\alpha= 
    \left[
\hspace{-12pt}
\begin{array}{ccc}
     & -a\sin\alpha\cos\theta\\
     &-a\sin\alpha\sin\theta\\
     &a\cos\alpha
\end{array}
\right]\; \;\;
 X_\theta= 
    \left[
\hspace{-12pt}
\begin{array}{ccc}
     & -(b+a\cos\alpha)\sin\theta\\
     &+(b+a\cos\alpha)\cos\theta\\
     &0
\end{array}
\right]\; \; \; 
    \left.
\hspace{-12pt}
\begin{array}{lll}
     & E=a^2 \\
     &F=0\\
     &G=(b+a\cos\alpha)^2
\end{array}
\right.
\end{align*}
\begin{align*}
    &ds^2 = a^2 d\alpha^2 + (b+a\cos\alpha)^2 d\theta^2 \; \; \;\; A = \int_{0}^{2\pi}\int_{0}^{2\pi}a(b+a\cos\alpha)\,d\alpha\,d\theta= 4\pi^2  ab
\end{align*}
